        \begin{exhibitbox}[主题证伪路径]
        \scriptsize
\resizebox{\exhibitwidth}{!}{
\begin{tabular}{llll}
\toprule
\textbf{风险} & \textbf{等级} & \textbf{触发} & \textbf{模型处理} \\
\midrule
WF6订单/价格证伪 & L4 & 公告继续否认大额长单或价格传闻 & 下调电子特气市场锚 \\
钨价快速回落 & L3 & APT/钨粉连续4周下行 & 下调钨股盈利和资源溢价 \\
制冷剂价差回落 & L3 & R32/R125/R134a价差跌回2025低位 & 下调氟化工PE \\
晶圆厂认证慢 & L3 & 客户导入或批量供货低于预期 & 直接产能折扣 \\
主题成交降温 & L4 & 龙头破20日均线且无公告支撑 & 等待财报而非追价 \\
\bottomrule
\end{tabular}
}
\end{exhibitbox}

\begin{exhibitbox}[重点标的催化与失效条件]
\scriptsize
\resizebox{\exhibitwidth}{!}{
\begin{tabular}{llll}
\toprule
\textbf{代码} & \textbf{名称} & \textbf{催化} & \textbf{失效条件} \\
\midrule
603379 & 三美股份 & R32/R125/R134a配额约束和价差修复、电子级HF/氟化液客户验证推进 & 制冷剂价格低于配额景气预期，或含氟材料/电子级产品认证低于预期 \\
002842 & 翔鹭钨业 & 钨矿/钨粉价格维持高位、出口许可偏紧、光伏钨丝与军工高温材料需求兑现 & 钨价快速回落、矿端供给释放或下游硬质合金需求承压导致Q2利润环比下滑 \\
300037 & 新宙邦 & R32/R125/R134a配额约束和价差修复、电子级HF/氟化液客户验证推进 & 制冷剂价格低于配额景气预期，或含氟材料/电子级产品认证低于预期 \\
605020 & 永和股份 & R32/R125/R134a配额约束和价差修复、电子级HF/氟化液客户验证推进 & 制冷剂价格低于配额景气预期，或含氟材料/电子级产品认证低于预期 \\
002378 & 章源钨业 & 钨矿/钨粉价格维持高位、出口许可偏紧、光伏钨丝与军工高温材料需求兑现 & 钨价快速回落、矿端供给释放或下游硬质合金需求承压导致Q2利润环比下滑 \\
600160 & 巨化股份 & R32/R125/R134a配额约束和价差修复、电子级HF/氟化液客户验证推进 & 制冷剂价格低于配额景气预期，或含氟材料/电子级产品认证低于预期 \\
002407 & 多氟多 & R32/R125/R134a配额约束和价差修复、电子级HF/氟化液客户验证推进 & 制冷剂价格低于配额景气预期，或含氟材料/电子级产品认证低于预期 \\
600378 & 昊华科技 & WF6/NF3/电子级氯化氢等品类价格和国产晶圆厂认证放量 & WF6价格/订单缺乏公告验证，或单一产品预期被证伪导致估值锚下移 \\
300346 & 南大光电 & WF6/NF3/电子级氯化氢等品类价格和国产晶圆厂认证放量 & WF6价格/订单缺乏公告验证，或单一产品预期被证伪导致估值锚下移 \\
000657 & 中钨高新 & 钨矿/钨粉价格维持高位、出口许可偏紧、光伏钨丝与军工高温材料需求兑现 & 钨价快速回落、矿端供给释放或下游硬质合金需求承压导致Q2利润环比下滑 \\
\bottomrule
\end{tabular}
}
\end{exhibitbox}

\section{后续监测清单}
第一，跟踪公司公告而不是题材传闻。WF6价格、规格、订单、客户和扩产只要不是公司正式披露，都只能作为情绪变量。第二，跟踪Q2/Q3财报中的收入、毛利率和存货变化；如果价格上涨没有进入毛利率，说明公司只是题材暴露而非利润暴露。第三，跟踪下游晶圆厂和存储扩产节奏；若先进存储资本开支下修，电子特气需求锚会被削弱。

第四，跟踪成交结构。若龙头日成交额维持高位且回撤能快速修复，市场锚可以保留；若成交额下降、主题扩散到低相关标的且公告验证缺位，应主动降低情绪权重。第五，跟踪政策端。钨出口许可、战略矿产目录、制冷剂配额和PFAS限制都可能改变供给约束和估值倍数。

\section{报告使用边界}
本报告适合用于建立产业链地图、筛选后续调研优先级和设定监测触发条件，不适合直接替代交易系统。最重要的结论是“分清直接受益和概念映射”：中船特气、中巨芯属于WF6直接线；巨化、三美、永和属于制冷剂和氟化工现金流线；厦门钨业、章源钨业属于资源价格线；金宏、凯美特、和远等需要更多产品和客户证据后再提高权重。
